\documentclass[11pt,a4paper]{article}
\usepackage[margin=24mm]{geometry}
\usepackage[T1]{fontenc}
\usepackage{lmodern}
\usepackage{amsmath,amssymb,graphicx,booktabs,tabularx,array,placeins}
\usepackage[font=small,labelfont=bf]{caption}
\usepackage{hyperref}
\hypersetup{hidelinks,pdftitle={Accelerating integral decomposition with Kira, machine learning, and exact reductions},pdfauthor={Rushil Saraswat}}
\setlength{\parindent}{1em}
\setlength{\parskip}{3pt}
\setlength{\emergencystretch}{2em}
\newcommand{\Q}{\mathbb Q}
\newcolumntype{Y}{>{\raggedright\arraybackslash}X}
\title{Accelerating integral decomposition with Kira, machine learning, and exact reductions}
\author{Rushil Saraswat}
\date{October 19, 2026}
\begin{document}
\maketitle
\input{paper_abstract}

\section{Introduction}
Feynman integrals arise from integration over internal loop momenta in higher-order perturbative corrections to scattering amplitudes and particle decay amplitudes. Many integrals in a calculation share propagators and differ only in their powers or numerators. Integration-by-parts (IBP) identities relate these integrals and express them in terms of a smaller set of master integrals~\cite{Weinzierl2022,ChetyrkinTkachov}. This decomposition supplies the coefficients needed to assemble an amplitude; the evaluation of the master integrals is a subsequent calculation~\cite{Weinzierl2022}.

Kira automates the generation and solution of the IBP system, and its current methods already remove unnecessary seeds and equations~\cite{KiraOriginal,Kira3}. The remaining cost depends on the integral structure and the requested targets. Generating fewer useful equations can reduce work, but a restricted system can also fail to reach the required basis and require a larger calculation. We therefore optimize complete reduction time, including prediction, algebra, and recovery. Seed-count reduction alone does not measure this objective.

Our method first identifies structures that admit an exact reduction formula. It applies those formulas to complete requests, individual targets, and recognized intermediate sectors. For the remaining supported native calculations, a trained cost model chooses a seed region before Kira generates the identities. Exact algebra determines every output coefficient. Related machine-assisted searches motivate the seed intervention~\cite{Song2025,TubeSeeding2026}, while analytic one-loop reduction motivates the formula stage~\cite{DennerDittmaier2006,Hu2025}. We derive the particular formulas and cost objective used here in the appendices.

The new extension treats two-point sectors with zero external invariant, including unequal, equal, and zero masses and polynomial numerators. It also retrains native seed selection on measured costs of target batches. We compare the integrated method with the preceding application and with a fixed seed choice. These controls distinguish improvements from exact identities and their integration from improvements attributable to ML. Earlier equation-selection and row-selection studies explain which other learned decisions we retain. The companion notes contain the full historical experiments and broader coverage audits.

\section{The decomposition problem and Kira}
\label{sec:physics}
We first define the notation used in this paper. An integral family specifies inverse propagators, masses, and external scalar products. In dimensional regularization, we write~\cite{Weinzierl2022,KiraOriginal}
\begin{equation}
 I_{\boldsymbol a}(\boldsymbol x,d)=\int\prod_{h=1}^{L}
 \frac{d^d\ell_h}{i\pi^{d/2}}\prod_{j=1}^{N}D_j^{-a_j},
 \qquad D_j=Q_j^2-m_j^2+i0,\qquad d=4-2\epsilon.
 \label{eq:integral}
\end{equation}
Here $L$ is the number of loops, $\ell_h$ is loop momentum $h$, $N$ is the number of integral coordinates, and $Q_j$ is a specified linear combination of loop and external momenta. The quantity $m_j$ is the mass on line $j$, $\boldsymbol x$ contains squared masses and external invariants, $\epsilon$ is the dimensional regulator, and $i0$ specifies the Feynman pole prescription. A positive integer $a_j$ gives a denominator power, zero removes that factor, and a negative integer gives a numerator power. Auxiliary inverse propagators complete the coordinates for loop scalar products when the physical lines do not suffice~\cite{KiraOriginal}.

A sector consists of integrals with the same positive-index positions. Removing a positive position gives a sub-sector. A seed $\boldsymbol\sigma$ is an index vector at which we generate identities~\cite{KiraOriginal}. For a loop or external momentum $v$, the IBP relation is~\cite{ChetyrkinTkachov}
\begin{equation}
 0=\int\prod_h\frac{d^d\ell_h}{i\pi^{d/2}}
 \frac{\partial}{\partial\ell_k^\mu}
 \left(v^\mu\prod_jD_j^{-\sigma_j}\right)
 =\sum_u A_{\boldsymbol\sigma,o;u}(\boldsymbol x,d)
 I_{\boldsymbol a^{(u)}}.
 \label{eq:identity}
\end{equation}
The operator label $o$ specifies $(k,v)$, and $A_{\boldsymbol\sigma,o;u}$ is the coefficient of integral $u$. Appendix~\ref{app:derivation} expands the derivative into index shifts and explains why the coefficients are rational. Each seed supplies several operators and therefore several equations.

Kira orders integrals, uses symmetries and other valid relations, selects equations, and performs elimination and back-substitution~\cite{KiraOriginal,Kira3,Laporta}. For a target $T$, the result has the form
\begin{equation}
 T=\sum_{m=1}^{n_M}r_{Tm}(\boldsymbol x,d)M_m,
 \qquad r_{Tm}\in\Q(\boldsymbol x,d).
 \label{eq:reduction}
\end{equation}
Here $M_m$ is master integral $m$, $n_M$ is the number of masters, and $\Q(\boldsymbol x,d)$ is the field of rational functions with rational coefficients. Exact row operations give Eq.~\eqref{eq:reduction}, as derived in Appendix~\ref{app:derivation}. A remaining free integral in a restricted system does not by itself establish a master of the full family. We check completion against a reference basis. Kira can solve symbolically or use finite fields and reconstruct rational functions~\cite{Kira2}; this study leaves its reconstruction algorithm unchanged.

The three learned interventions occur at distinct points (Table~\ref{tab:places}). A seed policy changes which identities are generated. An equation policy selects seed--operator pairs. A row policy chooses among allowed elimination steps after coefficients exist. All three can change cost without predicting the physical coefficients themselves.
\begin{table}[!htbp]\centering\small
\caption{Learned decisions studied here. The last column states their use in the measured method.}
\label{tab:places}
\begin{tabularx}{\linewidth}{@{}lYY@{}}\toprule
Decision & Work affected & Use in this study\\\midrule
Seed region & Generation and solution of the IBP system & Cost-trained policy in supported native requests\\
Equation retention & Construction of selected IBP coefficients & Separately validated native filter\\
Elimination row & Intermediate expressions and row operations & Omitted from the retained application\\\bottomrule
\end{tabularx}
\end{table}

\section{Exact reductions before native elimination}
\label{sec:exact}
The exact stage uses the propagator structure to determine whether a formula applies. For independent loop components, the product measure and separated integrand factorize Eq.~\eqref{eq:integral}; Appendix~\ref{app:derivation} proves the resulting product of component reductions. Identical components share their calculation. Connected graphs require their own identities, so a high loop count alone does not select this shortcut.

Table~\ref{tab:exact-domains} summarizes the retained domains. For generic one-loop sectors, we resolve numerator scalar products along the external directions and average the perpendicular components. This yields scalar integrals whose raised powers are lowered by an exact IBP recurrence. Appendix~\ref{app:tensor-reduction} derives both steps. Their coefficients are rational in $d$, and keeping known linear dimension factors separate reduces repeated symbolic simplification. A singular matrix can prevent those particular recurrences from being used, which motivates the new two-point rule.
\begin{table}[!htbp]\centering\small
\caption{Exact identities used by the application. Bounds describe the implementation and do not classify all reducible integrals.}
\label{tab:exact-domains}
\begin{tabularx}{\linewidth}{@{}YY@{}}\toprule
Integral structure & Identity and accepted domain\\\midrule
Generic one-loop sectors & Tensor angular average and scalar IBP recurrence; rational kinematics, at most five coordinates, numerator and excess positive powers at most six; required matrices invertible\\
Zero-invariant two-point sectors & New endpoint formula; equal, unequal, or zero masses; rational kinematics, at most five coordinates, numerator and excess positive powers at most four\\
Symbolic two-point masses & New scalar mass-derivative formula at $p^2=0$; individual powers at most six and total power at most eight\\
Connected one-scale trees & Successive massless bubble integration; massless two-point root or one massive vacuum root; implemented through thirteen loops\\
Independent products & Product of accepted component reductions; repeated components computed once\\\bottomrule
\end{tabularx}
\end{table}

\subsection{Two-point sectors at zero external invariant}
Let $D_1=\ell^2-x$ and $D_2=(\ell-p)^2-y$, where $x$ and $y$ are squared masses and $p^2=0$. Write $T(x)$ for the unit-power tadpole with squared mass $x$. Feynman parametrization gives
\begin{equation}
 B_{11}(x,y)=\frac{T(x)-T(y)}{x-y},\qquad
 B_{ab}(x,y)=\frac{\partial_x^{a-1}\partial_y^{b-1}B_{11}(x,y)}{(a-1)!(b-1)!}.
 \label{eq:light-main}
\end{equation}
Here $a$ and $b$ are positive propagator powers, and derivatives treat the two masses as independent before any specialization. Appendix~\ref{app:lightlike} derives this relation, its equal-mass limit, and the finite tensor formula. At equal masses, $B_{ab}(x,x)=T_{a+b}(x)$, where $T_n$ is the tadpole with denominator power $n$. A zero-mass tadpole vanishes in dimensional regularization~\cite{Weinzierl2022}.

For polynomial numerators, the same parameter shift leaves one radial denominator. Angular moments reduce the numerator to a finite sum, and a change from the Feynman parameter to the interpolated squared mass expresses each term through the two endpoint tadpoles. The result is a rational coefficient of $T(x)$ and $T(y)$, with no inverse external Gram matrix. This also applies when a generic one-loop recurrence reaches such a two-line sub-sector. Applying the identity at that intermediate step extends the complete reduction beyond requests that begin as bubbles.

The application tests the stated conditions before using a formula. It retains native reduction for unsupported targets, symbolic numerators outside the accepted domain, and recurrence limits. A request can contain both accepted and remaining targets; the application combines their exact outputs. These checks determine mathematical applicability. They do not ask a learned model to certify an identity.

\subsection{Connected trees and higher loop counts}
Successive massless bubble integrations also reduce specific connected trees to a unit-power integral of the same tree. For a massless two-point root, the coefficient is a rational function of $d$ and the external invariant. For a massive vacuum root, it is rational in $d$ and the one squared mass. Appendix~\ref{app:trees} derives these statements from Gamma-function products and their integer shifts. This family includes connected high-loop integrals, while independent products provide a separate high-loop test. Neither construction represents general coupled multiscale scattering.

\section{Learning the remaining Kira cost}
\label{sec:learning}
For native requests, the model selects a complete seed region from three allowed choices. Every choice uses exact Kira algebra and retains the same recovery checks. The actions differ in restrictions on the positive powers and numerator powers of seeds; Appendix~\ref{app:sparse} defines them. We measure the elapsed time of each action, including its unsuccessful attempts, so training favors choices that reduce the full cost.

The new training set contains 34 planar and nonplanar two-loop box batches, and a separate set of 16 batches selects the model. Each batch contains one to six targets with unit physical propagator powers and numerator rank one through four. Twenty-one features describe target counts, numerator distributions, and loop incidence in the physical propagators. Family names and the numerical kinematic profile are excluded from the features. Appendix~\ref{app:families-new} defines the integrals and split construction. Two additional training batches cover single-coordinate numerator distributions after a validation-only coverage check. No reserved timing results enter this revision. Whole target batches are disjoint across training, validation, and testing, although individual index vectors can recur.

We train cost-sensitive trees and bootstrap ensembles, using the tree constructions of Refs.~\cite{CART1984,Breiman2001Forest}. Each leaf stores the mean measured cost of each seed action. A split reduces the sum of chosen action costs; Appendix~\ref{app:shared} derives this objective and lists the candidate settings. Validation chooses both the model and a fixed control action. The new model is enabled only if validation supports a saving of at least one percent without more unverified results. A required predicted-saving margin limits changes near a tie. A separate learned failure score checks the proposed action before it runs. Predicted risky actions and inputs outside the training feature range use a conservative region that completed every training and validation request. The tests record each admission decision and any remaining failure.

This policy changes computational choices and leaves all coefficients to exact algebra. If the selected region is insufficient, the native procedure expands it and includes the extra time. A failed action can therefore erase a predicted saving. Our cost objective penalizes unverified actions, and the admission model learns from their observed outcomes. This reduces exposure to costly failed regions, while exact completion checks remain necessary for new requests. A failure label means that the bounded procedure did not return a verified complete reduction; it does not show that the integral cannot be reduced. The reserved tests count both completion and total time.

\subsection{Equation and row policies}
The equation-selection model ranks candidate seed--operator pairs before constructing their full coefficients. Its independently tested filter reduces complete time on its recorded workload. The current combined comparison keeps that filter off, so its benefit is reported separately. Multiplying the seed and equation speed ratios would assume an interaction that these measurements do not test.

The original next-row model scores allowed elimination steps in Kira's changing system. Its inference and selection costs exceeded the algebraic savings in the retained comparison. A later compiled cost policy reduced this overhead, but its interval still included equal time. We therefore omit learned row selection from the final application and retain Kira's exact elimination rules. Section~\ref{sec:results} reports the corresponding measurements.

\section{Comparison design}
\label{sec:design}
We freeze the formulas, trained model, fixed action, and reserved requests before measuring results. The new comparison contains 68 requests: 20 zero-invariant tensor bubbles, three bubbles with symbolic masses, 12 independent products at loop orders two through thirteen, three larger sectors that contain a zero-invariant two-line sub-sector, three ordinary one-loop tensor requests, three connected vacuum-tree requests, and 24 native box batches. Appendix~\ref{app:families-new} gives their kinematics and target construction. The connected vacuum requests have one, seven, and thirteen loops.

An independent enlarged-region Kira calculation supplies the coefficients for each one-loop component and native box request. We identify translation-equivalent equal-mass tadpoles before comparing coefficients. Product reductions retain their factored representation, and every component is checked. The vacuum requests check the unchanged exact tree path against its preceding result; they test preservation of that path, while the derivation and earlier Gamma-product checks establish its analytic basis. Development also checks the new identities against signed-index IBP residuals and tests symbolic mass expressions after independent rational substitutions.

For each reserved request, we run the integrated application, the same application with a fixed native seed choice, and the preceding published application. Each measurement starts a fresh process and includes imports, model loading, prediction, exact formulas, native generation and elimination where used, recovery, output, and exit. We use one computational thread per application, three repetitions, and rotating method order on the same worker. The external benchmark reference and its final coefficient comparison lie outside this timer. Native reference preparation supplies protected master indices where required; its coefficients are used only for the external check. Thus these native times measure reduction with the stated basis, rather than the cost of discovering that basis from scratch.

Each application has a 20-second bound. Completion requires agreement with the independent reference in all repetitions. Ratios use only requests that both compared methods complete; all remaining requests stay in the completion counts. We sum the per-request median times and form their ratio. Appendix~\ref{app:statistics} defines this estimator and its request-resampling interval. Timing uncertainty reflects this workload and hardware, not universal performance across integral families.

\input{lightlike_results}

\section{Conclusion}
\input{paper_conclusion}
\appendix
\section{IBP identities, row operations, and exactness}
\label{app:derivation}
The starting total-derivative identity follows from integration by parts in a region where the integral converges, followed by dimensional analytic continuation~\cite{ChetyrkinTkachov,Weinzierl2022}. Write $P_{\boldsymbol\sigma}=\prod_iD_i^{-\sigma_i}$ and $\partial_k=\partial/\partial\ell_k$. The product rule gives
\begin{align}
 \partial_k\!\cdot(vP_{\boldsymbol\sigma})
 &=(\partial_k\!\cdot v)P_{\boldsymbol\sigma}
 -\sum_i\sigma_i(v\!\cdot\partial_kD_i)
 D_i^{-\sigma_i-1}\prod_{j\ne i}D_j^{-\sigma_j}.
 \label{eq:product-rule}
\end{align}
For $v=\ell_h$, the first factor is $d\delta_{kh}$; for an external momentum it is zero. Completeness of the inverse-propagator coordinates permits us to write $v\!\cdot\partial_kD_i=u_i+\sum_jv_{ij}D_j$, where $u_i$ and $v_{ij}$ are coefficients fixed by the propagators and kinematics, with the invariant and mass terms included in $u_i$. Integrating Eq.~\eqref{eq:product-rule} therefore gives
\begin{equation}
 0=(\partial_k\!\cdot v)I_{\boldsymbol\sigma}
 -\sum_i\sigma_i u_i I_{\boldsymbol\sigma+\boldsymbol e_i}
 -\sum_{i,j}\sigma_i v_{ij}
 I_{\boldsymbol\sigma+\boldsymbol e_i-\boldsymbol e_j}.
 \label{eq:index-shifts-new}
\end{equation}
The unit vector $\boldsymbol e_i$ raises index $i$ by one. This derivation explains how a seed and an operator generate a row of the reduction system.

Let $q$ be a source row and $k$ a pivot column with $A_{qk}\ne0$. Solving the identity $\sum_jA_{qj}I_j=0$ for $I_k$ gives $I_k=-\sum_{j\ne k}(A_{qj}/A_{qk})I_j$. Subtracting $A_{rk}/A_{qk}$ times that zero identity from row $r$ yields
\begin{equation}
 A'_{rj}=A_{rj}-\frac{A_{rk}}{A_{qk}}A_{qj},\qquad A'_{rk}=0.
 \label{eq:row-new}
\end{equation}
Every step preserves the represented relation. Repetition until only prescribed master columns remain gives Eq.~\eqref{eq:reduction}. Removing seeds or equations preserves the validity of the retained identities, but can prevent this elimination from finishing. We therefore require reduction to the prescribed basis and exact agreement with the reference coefficients. Generic kinematics avoid singular pivot denominators; special kinematic limits require their own treatment.


For independent components, write the loop momenta as disjoint groups and the integrand as $\prod_b f_b$. In a convergence domain, repeated integration gives $I=\prod_b\int f_b$. Dimensional analytic continuation preserves that identity~\cite{Weinzierl2022}. Substituting each exact component reduction therefore gives
\begin{equation}
 I=\prod_b\left(\sum_m r_{bm}M_{bm}\right).
 \label{eq:product-reduction}
\end{equation}
Our product output retains this factored form. It avoids expanding the Cartesian product of master terms and uses the same representation in all timed methods.

\section{The seed regions and equation score}
\label{app:sparse}
For seed $\boldsymbol\sigma$, define excess positive power $D=\sum_i\max(\sigma_i-1,0)$, numerator rank $R=\sum_i\max(-\sigma_i,0)$, and positive sector $P=\{i:\sigma_i>0\}$. Let $h$ be the number of absent physical lines relative to the largest top sector, and let $D_T,R_T$ be the largest target values. The distance from the target numerator rectangles is defined by
\begin{equation}
 q(\boldsymbol\sigma,\mathcal T)=\min_{\boldsymbol a\in\mathcal T}
 \sum_{i\notin P}\max\{0,-\sigma_i-\max(0,-a_i)\}.
 \label{eq:seed-distance}
\end{equation}
Projection onto the interval from zero to each target numerator power leaves precisely this positive excess. Summing the distances and minimizing over targets gives Eq.~\eqref{eq:seed-distance}. We define the region by
\begin{align}
 q&\leq\max\{0,\lfloor\theta_0+\theta_1h+\theta_2D+\theta_6R\rfloor\},\nonumber\\
 R&\leq\max\{0,\lfloor R_T+\theta_3-\theta_4h\rfloor\},\qquad
 D\leq\max\{0,\lfloor D_T+\theta_5\rfloor\}.
 \label{eq:seed-regions}
\end{align}
These inequalities specify a search family rather than a sufficiency theorem. Native outer bounds, retained sector corners, and admissible target and master seeds also apply. The three choices use
\begin{align}
 \boldsymbol\theta_{\rm wide}&=(2.5,0,0,1.25,0,1,-0.25),\nonumber\\
 \boldsymbol\theta_{\rm tight}&=(1,0,0,0,1,1,0),\nonumber\\
 \boldsymbol\theta_{\rm medium}&=(1.5,0,0,0.5,0,1,0).
 \label{eq:seed-values}
\end{align}
The native procedure retains its subsector truncation and recovery rules. The model selects one complete choice before generating equations.

The separately tested equation classifier has seventeen features: four family and request dimensions; three seed counts; two differences from target power maxima; three sector descriptors; exact target membership and minimum index distance; the operator index; and the two target maxima. Positive labels mark equations retained by a complete reference reduction. Sampling at most 2,000 records of each class $y$ in batch $c$ gives each sampled record weight $N_{cy}/(n_{cy}N_c)$, where $N_{cy}$ and $n_{cy}$ are the native and sampled class counts and $N_c$ is the full batch count. Multiplying this weight by $n_{cy}$ preserves the class fraction $N_{cy}/N_c$. A leaf with weighted totals $W_0,W_1$ returns $p=W_1/(W_0+W_1)$, and a fixed threshold retains equations. The decision-tree construction follows Ref.~\cite{CART1984}. This empirical score does not certify sufficiency or probability calibration; the native completion check remains necessary.

\section{Measured-cost learning}
\label{app:shared}
Let $C_{ia}$ be the measured complete time of training request $i$ under action $a$. Unverified actions receive $C_{ia}\geq60$ seconds. This is a training penalty, and test failures are reported as failures. A leaf containing requests $S$ assigns action cost
\begin{equation}
 \widehat C_a(S)=\frac{1}{|S|}\sum_{i\in S}C_{ia},\qquad
 a_S=\arg\min_a\sum_{i\in S}C_{ia}.
 \label{eq:cost-leaf}
\end{equation}
The second formula follows because the first denominator does not depend on the action. A proposed split $S=S_1\cup S_2$ improves the empirical chosen cost by
\begin{equation}
 \Delta(S_1,S_2)=\min_a\sum_{i\in S}C_{ia}
 -\min_a\sum_{i\in S_1}C_{ia}
 -\min_a\sum_{i\in S_2}C_{ia}.
 \label{eq:cost-split}
\end{equation}
This is the parent cost minus the sum of child costs. We require $\Delta>0.002|S|$ seconds before splitting. This cost-sensitive objective uses the recursive tree and bootstrap-ensemble constructions of Refs.~\cite{CART1984,Breiman2001Forest}.

For each target, let $u$ and $v$ be its two nonnegative auxiliary numerator powers and let $R=u+v$. The 21 inputs are batch size; minimum, maximum, and mean $R$; number of distinct $R$ values; minimum, maximum, and mean of each of $u,v$; counts with $u=0$ and $v=0$; sums of $uv$ and $(u-v)^2$; number of physical propagators containing both loops; largest number of physical propagators incident on a loop; and counts with $R=1,2,3,4$. All are available before reduction.

We compare 1, 12, or 32 trees, maximum depths 1 through 5, minimum leaf populations 2, 4, or 8, and required fractional savings $\eta=0,0.03,0.08,0.15$. An ensemble averages leaf cost vectors. A single tree uses the complete training set; each ensemble tree samples training requests with replacement. The fixed action $b$ minimizes total validation cost among actions that completed every training and validation request. If $a_*$ has the least predicted cost, we choose it only when
\begin{equation}
 \widehat C_{a_*}+0.002 < (1-\eta)\widehat C_b.
 \label{eq:cost-margin}
\end{equation}
The allowance is two milliseconds per prediction during selection. Equation~\eqref{eq:cost-margin} is the defined minimum predicted saving, obtained by adding prediction overhead to the selected cost. Validation compares candidates by unverified count, complete cost, tree count, and depth, in that order. Deployment additionally requires at least one percent less validation cost and a change from the fixed action. We also fit a failure classifier for each action from its verified/unverified training labels. Each classifier has twelve trees of depth at most three and at least two requests per child. One tree uses all requests, and eleven use bootstrap samples. A leaf returns its failed fraction $p_f$. Binary Gini impurity for $n$ requests with $k$ failures is
\begin{equation}
 n\left[1-(k/n)^2-(1-k/n)^2\right]=\frac{2k(n-k)}{n}.
 \label{eq:failure-gini}
\end{equation}
The equality follows by expansion. Splits maximize the decrease in this quantity, following the classification-tree construction of Ref.~\cite{CART1984}. Averaging the leaf fractions gives an empirical risk score for each action, without a claim of calibrated failure probability. Validation selects a threshold from $0,0.05,0.15,0.30$ jointly with the cost policy. If the proposed score exceeds that threshold, or any feature lies outside its training range, the policy selects the fully verified fixed region. Validation requires no unverified selected actions before enabling the model. This guard uses observed evidence; a zero score is not a proof that an unseen reduction will complete. The application timer includes both predictions and any subsequent recovery.

\section{Physical test inputs}
\label{app:families-new}
All profiles below define squared quantities in a common fixed unit. The zero-invariant families use shifts $(0,p,q,r)$ and masses $(x,y,0,0)$, with
\begin{align}
 &p^2=0,\quad p\cdot q=-h-1,\quad p\cdot r=-h-2,\nonumber\\
 &q^2=-2h-7,\quad q\cdot r=-h-3,\quad r^2=-3h-13.
 \label{eq:light-input}
\end{align}
The unequal case has $(x,y)=(h+3,h+11)$, the equal case $(h+3,h+3)$, the mixed case $(h+3,0)$, and the massless case $(0,0)$. Targets have the form $(1+\delta_1,1+\delta_2,-b_1,-b_2)$ with $\delta_1+\delta_2=\delta$ and $b_1+b_2=R$. For class position $j=0,1,2,3$, we use $R=0,\ldots,4$, $\delta=\min(R,3)$, and $h=211+17j+R$. Each batch contains at most three distinct assignments. These definitions give the 20 tensor requests.

The three symbolic requests retain the same zero-invariant geometry with $h=301,302,303$ and squared masses $(x,y)$, $(x,x)$, and $(x,0)$ as symbols. Their positive-power pairs are $(1,1),(2,1),(1,2),(2,2)$, with zero numerator powers. Independent products contain $L$ identical components at $h=351+L$, $R=2$, and $\delta=1$, with unequal masses for odd $L$ and equal masses for even $L$. Each target is a product of the same component assignment, and $L=2,\ldots,13$. No factor connects different loop momenta.

The embedded-sector requests promote the $q$ line to a physical denominator with squared mass $421+j$ and use $h=391+j$ in the three nonzero-mass classes, $j=0,1,2$. Their targets are $(2,1,1,-1)$, $(1,2,1,-1)$, and $(2,1,0,-1)$. The first two require a three-line reduction, and the last explicitly tests its zero-invariant sub-sector.

The ordinary tensor checks have $P=2,3,4$ physical lines, one auxiliary coordinate, shifts $q_1=0$, $q_{i+1}=p_i$, and profile $h=433+P$. Their masses and scalar products are
\begin{equation}
 m_i^2=h+11i^2\ (i\leq P),\quad m_{P+1}^2=0,\quad
 p_i^2=-(10h+17i),\quad p_i\cdot p_j=-(h+i+2j-2)\ (i<j).
 \label{eq:ordinary-input}
\end{equation}
Targets have numerator rank two and one excess positive power. The unchanged connected-vacuum checks use chains at $L=1,7,13$, the family of Appendix~\ref{app:trees}, and four target assignments per request. The input records retain the exact integer assignments for every batch in this finite comparison.

For the native planar box, let $p_1,p_2,p_3$ be massless, $2p_1\cdot p_2=s$, $2p_1\cdot p_3=t$, and $2p_2\cdot p_3=-s-t$. The seven physical line momenta are
\begin{equation}
 (\ell_1,\ell_1-p_1,\ell_1-p_1-p_2,\ell_2,
 \ell_2-p_1-p_2,\ell_2-p_1-p_2-p_3,\ell_1-\ell_2),
 \label{eq:planar-input}
\end{equation}
and the auxiliary momenta are $(\ell_1-p_3,\ell_2-p_1)$. The nonplanar routing follows Ref.~\cite{KiraOriginal}:
\begin{align}
 (Q_1,\ldots,Q_9)=(&\ell_1,\ell_2,p_1-\ell_1,p_2-\ell_2,
 p_1+p_2-p_3-\ell_1,\nonumber\\
 &p_1-\ell_1+\ell_2,p_1-p_3-\ell_1+\ell_2,
 \ell_1-p_3,\ell_2-p_1-p_2).
 \label{eq:nonplanar-input}
\end{align}
The first seven entries are physical and all masses vanish. Its convention gives $p_i^2=(p_1+p_2-p_3)^2=0$, $s=(p_1+p_2)^2$, and $t=(p_1-p_3)^2$. Expanding the squares gives $2p_1\cdot p_3=-t$ and $2p_2\cdot p_3=s+t$. The planar profiles use $(s,t)=(-1009-12h,-1021-14h)$ and the nonplanar profiles use $(-113-18h,-127-10h)$.

Every native target is $(1,1,1,1,1,1,1,-u,-v)$ with $u,v\geq0$ and $1\leq u+v\leq4$. Distinct batches contain one to six targets. Training uses 16 batches per graph at $h=19079$ and one additional batch per graph at $h=19083$ with numerator powers $(u,v)=(0,2),(0,3),(0,4)$, validation uses eight per graph at $h=19121$, and testing uses 12 per graph at $h=19139$. We exclude whole batches used by the preceding native and tensor studies. The same graph types occur in all three new splits, so the test measures new batches and kinematic points within these structures rather than transfer to an unseen topology.

\section{Time ratios and uncertainty}
\label{app:statistics}
Let $T_{imr}$ be complete elapsed time for request $i$, method $m$, and repetition $r$. For the set $\mathcal C$ completed by both methods in all repetitions, define
\begin{equation}
 S_m=\sum_{i\in\mathcal C}\operatorname{median}_{r=1,2,3}T_{imr},\qquad
 R=\frac{S_{\rm new}}{S_{\rm control}},\qquad
 \text{fractional saving}=1-R.
 \label{eq:time-ratio}
\end{equation}
These definitions compare identical completed requests. Dividing both times by the control time gives $R$, and subtracting gives the fractional saving. Completion counts remain separate because an unfinished reduction has no verified finite speedup. To estimate variation across requests, we draw $|\mathcal C|$ paired request indices with replacement and recompute $R$. We report the central 95\% of 2,000 such ratios, using the bootstrap construction of Ref.~\cite{Efron1979}. This interval describes the sampled request collection and does not establish performance for every integral family.

\section{Exact reduction of polynomial numerators}
\label{app:tensor-reduction}
\subsection{Resolving the loop momentum}
One-loop tensor reduction expresses loop-momentum numerators through scalar integrals. Established methods include tensor decompositions and generating functions for arbitrary numerator and denominator powers~\cite{DennerDittmaier2006,Hu2025}. We derive the particular identities used in our extension here. The algebra uses rational kinematic values and keeps $d$ symbolic. It applies to a one-loop sector and to independent products of such sectors. Applying it inside a coupled graph would require coefficients that depend on the other loop momenta, which this implementation does not support.

Write $D_i=(\ell-q_i)^2-m_i^2$ and let $S=\{i:a_i>0\}$ contain the denominator positions. The $q_i$ in this appendix are fixed external shifts; the infinitesimal prescription is understood. Choose $r\in S$, set $k=\ell-q_r$, and define $R_i=q_i-q_r$. The vectors $R_i$ with $i\in S\setminus\{r\}$ span the external directions seen by the denominators. Their Gram matrix and their contractions with $k$ are
\begin{equation}
 G_{ij}=R_i\cdot R_j,\qquad
 u_i=k\cdot R_i=\frac{D_r-D_i+m_r^2-m_i^2+R_i^2}{2}.
 \label{eq:tensor-u}
\end{equation}
The second equality follows by expanding $D_i=(k-R_i)^2-m_i^2$ and subtracting it from $D_r=k^2-m_r^2$. Assume that $G$ is invertible, and split $k$ into its projection onto this span and its perpendicular part:
\begin{equation}
 k_\parallel^\mu=\sum_{i,j\in S\setminus\{r\}}R_i^\mu(G^{-1})_{ij}u_j,
 \qquad k_\perp=k-k_\parallel.
 \label{eq:tensor-projection}
\end{equation}
Contracting with each $R_i$ gives $k_\perp\cdot R_i=0$. Contracting the parallel part with itself gives $k_\parallel^2=u^TG^{-1}u$, so
\begin{equation}
 T(D):=k_\perp^2=D_r+m_r^2-u^TG^{-1}u.
 \label{eq:tensor-radius}
\end{equation}
Empty matrices and sums give zero when $S$ has one element.

An inverse propagator $D_j$ with $a_j<0$ appears in the numerator and can be resolved into the same components. Define $v_{ji}=R_j\cdot R_i$, $R_{j,\perp}=R_j-\sum_{i,h}R_i(G^{-1})_{ih}v_{jh}$, and $z_j=k_\perp\cdot R_{j,\perp}$. Direct substitution gives
\begin{align}
 D_j&=B_j(D)-2z_j,\\
 B_j(D)&=D_r+m_r^2+R_j^2-m_j^2-2v_j^TG^{-1}u,\\
 H_{jh}&:=R_{j,\perp}\cdot R_{h,\perp}
       =R_j\cdot R_h-v_j^TG^{-1}v_h.
 \label{eq:tensor-numerator}
\end{align}
Thus the remaining task is to average products of the $z_j$ over perpendicular directions. All positive-power denominators depend on $k_\perp$ only through $T(D)$, so this angular average leaves the denominator unchanged.

\subsection{Angular moments and the finite numerator expansion}
Let $\nu=d-|S|+1$ be the dimension perpendicular to the denominator span. Rotational symmetry makes odd angular moments vanish. For $2h$ perpendicular components, the invariant symmetric tensor is the sum of all pairings of metric tensors, multiplied by a coefficient $C_h T^h$. Tracing two indices gives
\begin{equation}
 C_h(\nu+2h-2)=C_{h-1},\qquad C_0=1,
 \qquad C_h=\frac{1}{\prod_{b=0}^{h-1}(\nu+2b)}.
 \label{eq:tensor-trace}
\end{equation}
The trace factor contains $\nu$ from the pairing of the traced indices and $2h-2$ from their pairings with the other indices. Contracting the tensor with the external perpendicular vectors therefore gives
\begin{equation}
 \left\langle\prod_{t=1}^{2h}z_{j_t}\right\rangle
 =\frac{T(D)^h}{\prod_{b=0}^{h-1}(\nu+2b)}
 \sum_{\mathcal P}\prod_{(s,t)\in\mathcal P}H_{j_sj_t}.
 \label{eq:tensor-moment}
\end{equation}
Here $\mathcal P$ runs over the distinct pairings of the $2h$ positions, including their multiplicities when an index repeats. In particular,
\begin{equation}
 \langle z_i z_j\rangle=\frac{T H_{ij}}{\nu},\qquad
 \langle z_i z_j z_k z_l\rangle
 =\frac{T^2(H_{ij}H_{kl}+H_{ik}H_{jl}+H_{il}H_{jk})}{\nu(\nu+2)}.
 \label{eq:tensor-low-moments}
\end{equation}
One can first derive these relations after a Wick rotation in an integer dimension, where the angular average is well defined, and then continue the rational tensor coefficients in $d$. This uses the dimensional-regularization convention of Eq.~\eqref{eq:integral}~\cite{Weinzierl2022}. Poles in $d$ remain exact rational poles.

For $b_j=-a_j>0$, expand $\prod_{j:a_j<0}(B_j-2z_j)^{b_j}$ and insert Eq.~\eqref{eq:tensor-moment}. The result is a finite polynomial in the positive-sector denominators:
\begin{equation}
 \left\langle\prod_{j:a_j<0}D_j^{b_j}\right\rangle
 =\sum_{\boldsymbol t}c_{\boldsymbol t}(d)\prod_{i\in S}D_i^{t_i},
 \qquad
 I_{\boldsymbol a}=\sum_{\boldsymbol t}c_{\boldsymbol t}(d)
 I_{\boldsymbol a_S-\boldsymbol t}.
 \label{eq:tensor-expansion}
\end{equation}
The vector $\boldsymbol a_S$ equals $\boldsymbol a$ on $S$ and vanishes elsewhere; the $t_i$ are nonnegative and vanish outside $S$. The coefficients also depend rationally on the fixed kinematics. Each power $D_i^{t_i}$ cancels the corresponding denominator power. If a cancellation creates a new negative index, that line leaves the positive sector. Repeating the angular step therefore decreases the positive-sector size whenever a numerator remains. A sector with no denominators is scaleless and vanishes in dimensional regularization~\cite{Weinzierl2022}. The process terminates in scalar sectors with nonnegative indices.

\subsection{Lowering powers before removing the numerator}
A second exact route starts with IBP. Define $A=\sum_j a_j$ and $Y_{ij}=m_i^2+m_j^2-(q_i-q_j)^2$. Expanding the total derivative with vector $\ell-q_i$ gives
\begin{align}
 0&=d I_{\boldsymbol a}
 -\sum_j a_j\int\frac{d^d\ell}{i\pi^{d/2}}
 \frac{2(\ell-q_i)\cdot(\ell-q_j)}{D_j}\prod_hD_h^{-a_h},\\
 2(\ell-q_i)\cdot(\ell-q_j)&=D_i+D_j+Y_{ij},\\
 0&=(d-A)I_{\boldsymbol a}
 -\sum_j a_j I_{\boldsymbol a+\boldsymbol e_j-\boldsymbol e_i}
 -\sum_j a_jY_{ij}I_{\boldsymbol a+\boldsymbol e_j}.
 \label{eq:tensor-ibp}
\end{align}
Here $\boldsymbol e_j$ is the unit index vector. This derivation holds for signed integer indices, so it includes numerator terms. Let $N_- =\{j:a_j<0\}$. Moving the negative-index terms to the right gives, for $i\in S$,
\begin{align}
 \sum_{j\in S}a_jY_{ij}I_{\boldsymbol a+\boldsymbol e_j}&=F_i(\boldsymbol a),\\
 F_i(\boldsymbol a)&=(d-A)I_{\boldsymbol a}
 -\sum_{j:a_j\ne0}a_j I_{\boldsymbol a+\boldsymbol e_j-\boldsymbol e_i}
 -\sum_{j\in N_-}a_jY_{ij}I_{\boldsymbol a+\boldsymbol e_j}.
 \label{eq:tensor-signed-system}
\end{align}
For target $\boldsymbol b$ with $b_k>1$, choose $\boldsymbol a=\boldsymbol b-\boldsymbol e_k$. If $Y_S$ is invertible, multiply this system by its inverse to obtain
\begin{equation}
 I_{\boldsymbol b}=\frac{1}{a_k}
 \sum_{i\in S}(Y_S^{-1})_{ki}F_i(\boldsymbol a).
 \label{eq:tensor-signed-recurrence}
\end{equation}
Every integral on the right has a smaller sum of positive indices than $\boldsymbol b$, and the numerator rank does not increase. When the positive powers reach one, the angular step removes the remaining numerator. With no negative indices, Eq.~\eqref{eq:tensor-signed-recurrence} reduces to the scalar Cayley recurrence derived earlier. These identities give a finite reduction to unit-power scalar sectors. They produce a spanning set; special kinematics can introduce further relations between its elements.

\subsection{Keeping the dimension poles in factored form}
The preceding identities restrict the form of every dimension-dependent denominator. The inverse Gram and Cayley matrices have rational entries at the chosen kinematic point, and Eq.~\eqref{eq:tensor-signed-recurrence} introduces only polynomial dependence on $d$. Equation~\eqref{eq:tensor-moment} introduces products of linear factors $d+j$, where $j$ is an integer. By induction through the recurrence, each coefficient therefore has the form
\begin{equation}
 r(d)=\frac{P(d)}{\prod_j(d+j)^{n_j}},\qquad
 P(d)\in\Q[d],\quad n_j\in\mathbb Z_{\geq0}.
 \label{eq:tensor-factored-field}
\end{equation}
The finite products contain only factors introduced by the angular moments. For addition, take $Q$ to be the least common multiple of two such denominators $Q_1$ and $Q_2$. Multiplying each summand by its missing factors gives
\begin{equation}
 \frac{P_1}{Q_1}+\frac{P_2}{Q_2}
 =\frac{P_1(Q/Q_1)+P_2(Q/Q_2)}{Q}.
 \label{eq:tensor-factored-add}
\end{equation}
The exponent of $d+j$ in $Q$ is the larger of its two exponents. Multiplication adds these exponents. A factor $d+j$ cancels exactly when $P(-j)=0$, because division with remainder gives $P(d)=(d+j)\widetilde P(d)+P(-j)$. We repeat this cancellation for each retained factor. These operations preserve the exact rational function and require only rational polynomial coefficients and integer factor multiplicities. They avoid general rational-function simplification in the new path. The earlier arithmetic checks compare these coefficients with a general rational-function implementation and native Kira reductions. The current application retains the fixed numerator-first recurrence; its new learned decision concerns native seed regions.

\section{Exact reduction at a zero two-point invariant}
\label{app:lightlike}
A singular Gram matrix prevents the projection used by the generic tensor recurrence. For a two-point sector with $p^2=0$, Feynman parameters give a direct reduction to tadpoles. Degenerate one-loop kinematics require such separate identities~\cite{DennerDittmaier2006}. We derive the identity used in this study below. It applies to a polynomial numerator and two positive denominator powers; it does not assume that the two masses differ.

\subsection{Scalar integrals and mass derivatives}
Let $x$ and $y$ denote squared masses, and define
\begin{equation}
 T(x)=\int\frac{d^d\ell}{i\pi^{d/2}}\frac{1}{\ell^2-x+i0},\qquad
 B_{ab}(x,y)=\left.\int\frac{d^d\ell}{i\pi^{d/2}}
 \frac{1}{(\ell^2-x+i0)^a[(\ell-p)^2-y+i0]^b}\right|_{p^2=0}.
 \label{eq:light-def}
\end{equation}
Here $a,b$ are positive integers. Dimensional regularization sets $T(0)=0$. With the normalization in Eq.~\eqref{eq:integral}, Wick rotation and the radial beta integral give~\cite{Weinzierl2022}
\begin{equation}
 T(x)=-\Gamma(1-d/2)(x-i0)^{d/2-1}.
 \label{eq:light-tad}
\end{equation}
For $a=b=1$, combine the denominators with a parameter $z$ and shift $\ell=k+(1-z)p$. The combined squared mass is $M^2(z)=zx+(1-z)y$, because $p^2=0$. Integrating over $k$ and then $z$ gives
\begin{align}
 B_{11}(x,y)
 &=\Gamma(2-d/2)\int_0^1 [zx+(1-z)y-i0]^{d/2-2}\,dz\nonumber\\
 &=\frac{T(x)-T(y)}{x-y},\qquad x\ne y.
 \label{eq:light-difference}
\end{align}
The second equality follows from the elementary antiderivative of a power and $\Gamma(2-d/2)=(1-d/2)\Gamma(1-d/2)$. Differentiating a propagator with respect to its squared mass raises its power. Therefore
\begin{equation}
 B_{ab}(x,y)=\frac{\partial_x^{a-1}\partial_y^{b-1}B_{11}(x,y)}{(a-1)!(b-1)!}.
 \label{eq:light-derivative}
\end{equation}
This operation preserves a rational combination of $T(x)$ and $T(y)$ because $\partial_xT(x)=(d/2-1)T(x)/x$. We take the derivatives before imposing relations between the masses. If $x=y\ne0$, the combined mass is independent of $z$, and the beta integral instead gives
\begin{equation}
 B_{ab}(x,x)=T_{a+b}(x),\qquad
 T_n(x)=\frac{\prod_{j=1}^{n-1}(d/2-j)}{(n-1)!\,x^{n-1}}T(x).
 \label{eq:light-equal}
\end{equation}
The formula for $T_n$ also follows by differentiating Eq.~\eqref{eq:light-tad}. These identities support independent symbolic squared masses as well as rational values. A relation involving distinct symbolic masses is a generic rational identity; equal masses use Eq.~\eqref{eq:light-equal} directly.

\subsection{Polynomial numerators}
Let $P(\ell)$ be a polynomial numerator. Combining the two denominators gives
\begin{equation}
 B_{ab}[P]=\frac{\Gamma(A)}{\Gamma(a)\Gamma(b)}\int_0^1 dz\,
 z^{a-1}(1-z)^{b-1}\int\frac{d^dk}{i\pi^{d/2}}
 \frac{P(k+(1-z)p)}{[k^2-M^2(z)+i0]^A},\quad A=a+b.
 \label{eq:light-parameter}
\end{equation}
For an auxiliary inverse propagator $(\ell-q_j)^2-m_j^2$, put $w_j=(1-z)p-q_j$. Its numerator factor becomes $k^2+2k\cdot w_j+w_j^2-m_j^2$. Both $w_i\cdot w_j$ and $w_j^2-m_j^2$ are polynomials of degree at most one in $z$ when $p^2=0$.

Rotational covariance gives zero for odd moments of $k$. Tracing two indices of an even moment determines its normalization recursively: if its pairing coefficient is $C_h$, then $C_h(d+2h-2)=C_{h-1}$ and $C_0=1$. Consequently,
\begin{equation}
 \left\langle\prod_{j=1}^{2h}(k\cdot w_j)\right\rangle
 =\frac{(k^2)^h}{\prod_{r=0}^{h-1}(d+2r)}
 \sum_{\mathcal P}\prod_{(i,j)\in\mathcal P}w_i\cdot w_j,
 \label{eq:light-angular}
\end{equation}
where $\mathcal P$ runs over pairings of the $2h$ positions. The trace recurrence follows because one pairing contributes $d$ and the remaining contractions contribute $2h-2$. This also derives the full-dimensional form of the angular identity used in Appendix~\ref{app:tensor-reduction}.

After the finite numerator expansion and angular average, define $c_{ht}(d)$ by
\begin{equation}
 \langle P(k+(1-z)p)\rangle=\sum_{h,t}c_{ht}(d)(k^2)^h z^t.
 \label{eq:light-polynomial}
\end{equation}
These coefficients also depend on the fixed external scalar products and numerator masses. They contain no loop momentum. To integrate each term, Wick rotate the radial integral, set $u=k_E^2/M^2$, and use
$\int_0^\infty u^{\alpha-1}(1+u)^{-A}du=\Gamma(\alpha)\Gamma(A-\alpha)/\Gamma(A)$ with $\alpha=d/2+h$. This gives
\begin{equation}
 J_{Ah}(M^2)=\int\frac{d^dk}{i\pi^{d/2}}\frac{(k^2)^h}{(k^2-M^2+i0)^A}
 =(-1)^{A-h}\frac{\Gamma(d/2+h)\Gamma(A-h-d/2)}{\Gamma(d/2)\Gamma(A)}
 (M^2-i0)^{d/2+h-A}.
 \label{eq:light-radial}
\end{equation}
The initial convergence region defines the integrals; analytic continuation gives the dimensional-regularization identity outside that region~\cite{Weinzierl2022}.

For $x\ne y$, set $\Delta=x-y$, $\rho=h-A+1$, and $\lambda=d/2-1$. Expand $(1-z)^{b-1}$ and then substitute $u=y+\Delta z$. For an integer $k\ge0$,
\begin{equation}
 \int_0^1 z^k(y+\Delta z)^{\lambda+\rho}\,dz
 =\frac{1}{\Delta^{k+1}}\sum_{j=0}^k\binom{k}{j}(-y)^{k-j}
 \frac{x^{\lambda+\rho+j+1}-y^{\lambda+\rho+j+1}}{\lambda+\rho+j+1}.
 \label{eq:light-endpoints}
\end{equation}
This follows by expanding $(u-y)^k$ and integrating each power of $u$. Dividing by the tadpole normalization in Eq.~\eqref{eq:light-tad}, define
\begin{equation}
 R_h=\frac{(-1)^{A-h+1}(d/2)_h(1-d/2)_{A-h-1}}{(a-1)!(b-1)!},
 \qquad (z)_n=\frac{\Gamma(z+n)}{\Gamma(z)}.
 \label{eq:light-factor}
\end{equation}
For a negative integer $n$, the last ratio means the inverse finite product $1/\prod_{r=1}^{-n}(z-r)$. Equations~\eqref{eq:light-parameter}--\eqref{eq:light-endpoints} now give the explicit rational reduction
\begin{align}
 B_{ab}[P]&=\sum_{h,t}c_{ht}(d)R_h\sum_{v=0}^{b-1}(-1)^v\binom{b-1}{v}
 \nonumber\\ &\quad\times\sum_{j=0}^{k}\binom{k}{j}\frac{(-y)^{k-j}}{\Delta^{k+1}}
 \frac{x^{\rho+j+1}T(x)-y^{\rho+j+1}T(y)}{\lambda+\rho+j+1},\nonumber\\[-2pt]
 &\quad k=a-1+t+v.
 \label{eq:light-final}
\end{align}
If one mass is zero, its endpoint vanishes in a convergence region and hence by analytic continuation; the implementation omits that scaleless endpoint before evaluating powers. If both masses vanish, the sector is scaleless. For equal nonzero masses, use the beta integral directly:
\begin{equation}
 B_{ab}[P]=\sum_{h,t}c_{ht}(d)R_h\,
 \frac{\Gamma(a+t)\Gamma(b)}{\Gamma(A+t)}x^{h-A+1}T(x).
 \label{eq:light-equal-tensor}
\end{equation}
Thus the equal-mass rule has no division by $x-y$. A sector with one positive denominator uses Eq.~\eqref{eq:light-radial} and the same angular expansion without a Feynman parameter.

At rational kinematics, the coefficients have polynomial numerators in $d$ and products of integer-shift linear factors $d+j$ in their denominators. The finite products in Eqs.~\eqref{eq:light-angular} and \eqref{eq:light-factor}, and the linear denominator in Eq.~\eqref{eq:light-endpoints}, prove this statement. We therefore use the exact factored arithmetic of Appendix~\ref{app:tensor-reduction}. The implemented tensor domain has at most five coordinates, numerator rank at most four, and at most four extra positive powers. Symbolic masses currently use the scalar derivative identity. Inputs outside these conditions continue through the existing exact rules and Kira.

\section{Connected one-scale tree reductions}
\label{app:trees}
We derive the previously retained tree rules to state precisely which high-loop structures they cover. Let $\mathcal T$ be a rooted tree with loop momenta $\ell_0,\ldots,\ell_{L-1}$ and parent $\pi(j)$ for every nonroot vertex. A massless two-point family is
\begin{equation}
 J_{\mathcal T}=\int\prod_{j=0}^{L-1}\frac{d^d\ell_j}{i\pi^{d/2}}
 \frac{1}{(\ell_0^2)^{a_0}[(\ell_0-p)^2]^c}
 \prod_{j=1}^{L-1}\frac{1}{(\ell_j^2)^{a_j}
 [ (\ell_j-\ell_{\pi(j)})^2]^{v_j}},\qquad p^2=s\ne0.
 \label{eq:tree-family}
\end{equation}
The massive vacuum version replaces the two root factors by $(\ell_0^2-m^2)^{-c}$ and has no external momentum. Each edge power $v_j$, root power $c$, and leaf self power $a_j$ is positive. Internal self powers can be nonpositive. Auxiliary powers vanish. The implementation admits $L\leq13$ and absolute physical powers at most sixteen. The derivation itself applies to finite trees with meromorphically defined integrals.

Temporarily use negative massless denominators. Feynman parametrization of a massless bubble, translation of its loop momentum, and a radial beta integral give
\begin{align}
 G(u,v)&=\frac{\Gamma(u+v-d/2)}{\Gamma(u)\Gamma(v)}
 \int_0^1z^{d/2-v-1}(1-z)^{d/2-u-1}\,dz\nonumber\\
 &=\frac{\Gamma(u+v-d/2)\Gamma(d/2-u)\Gamma(d/2-v)}
 {\Gamma(u)\Gamma(v)\Gamma(d-u-v)}.
 \label{eq:tree-bubble}
\end{align}
The remaining scale factor is $(-P^2-i0)^{d/2-u-v}$ for bubble momentum $P$. The beta identity used here follows by the same radial and parameter integrations as Appendix~\ref{app:lightlike}; analytic continuation supplies the dimensionally regularized relation~\cite{Weinzierl2022}. Recursive subintegral integration has established applications to massless propagator integrals~\cite{Ruijl2017Forcer}.

For each nonroot vertex, define
\begin{equation}
 u_j=a_j+\sum_{h:\pi(h)=j}\sigma_h,\qquad
 \sigma_j=u_j+v_j-d/2.
 \label{eq:tree-transfer}
\end{equation}
Integrating a leaf gives $G(u_j,v_j)$ and transfers power $\sigma_j$ to its parent. Repeating from the leaves to the root proves the product rule by induction. In the two-point case, set $\alpha=a_0+\sum_{h:\pi(h)=0}\sigma_h$; its dimensionless coefficient is
\begin{equation}
 C_{\mathcal T}=G(\alpha,c)\prod_{j=1}^{L-1}G(u_j,v_j).
 \label{eq:tree-coefficient}
\end{equation}
The total physical power $A=\sum_j a_j+\sum_{j>0}v_j+c$ gives scale power $(-s-i0)^{Ld/2-A}$. Let $M_{\mathcal T}$ set all physical powers to one and let $C_0$ be its Gamma product. Restoring the original propagator signs and dividing by this master gives
\begin{equation}
 J_{\mathcal T}=s^{-(A-2L)}\frac{C_{\mathcal T}}{C_0}M_{\mathcal T}.
 \label{eq:tree-ratio}
\end{equation}
The sign $(-1)^{A-2L}$ from the propagator convention cancels that from $(-s)^{-(A-2L)}$.

For a massive vacuum root, put $\sigma_0=\sum_{h:\pi(h)=0}\sigma_h$. Wick rotation and the substitution $x=m^2t$ give
\begin{equation}
 \frac{1}{\Gamma(d/2)}\int_0^\infty
 \frac{x^{d/2-\sigma_0-1}}{(x+m^2)^c}\,dx
 =(m^2)^{d/2-\sigma_0-c}
 \frac{\Gamma(d/2-\sigma_0)\Gamma(\sigma_0+c-d/2)}
 {\Gamma(d/2)\Gamma(c)}.
 \label{eq:tree-vacuum}
\end{equation}
Multiply the dimensionless Gamma ratio by the nonroot bubble factors to define $C_{\mathcal T}^{V}$. There are $P=2L-1$ physical lines and $A=c+\sum_{j>0}(a_j+v_j)$. Dividing by the unit-power vacuum master $M_{\mathcal T}^{V}$ yields
\begin{equation}
 V_{\mathcal T}=(-m^2)^{-(A-P)}
 \frac{C_{\mathcal T}^{V}}{C_0^{V}}M_{\mathcal T}^{V}.
 \label{eq:tree-vacuum-ratio}
\end{equation}
In either family, a subtree with $n_j$ vertices has $u_j=U_j-(n_j-1)d/2$ with integer $U_j$. This follows from Eq.~\eqref{eq:tree-transfer} by induction. Target and master Gamma arguments therefore have equal slopes in $d$ and integer differences in their constant terms. Repeated use of $\Gamma(z+1)=z\Gamma(z)$ gives
\begin{equation}
 \frac{\Gamma(z+n)}{\Gamma(z)}=
 \begin{cases}\prod_{k=0}^{n-1}(z+k),&n>0,\\
 1,&n=0,\\
 [\prod_{k=1}^{-n}(z-k)]^{-1},&n<0.
 \end{cases}
 \label{eq:tree-gamma-ratio}
\end{equation}
Thus the target-to-master coefficients are rational in $d$ and the one scale. Cancelling equal Gamma factors before expanding these finite products preserves the identity and avoids repeated symbolic work. This proves the reductions for the stated trees; additional coupled propagators or auxiliary numerators require another reduction.

\bibliographystyle{unsrt}
\bibliography{paper_references}
\end{document}
